\subsection{Molecular simulation details}

To investigate the diffusive dynamics of the atoms in realistic grains present in experimental settings,\cite{thiyaneshwaran2018nucleation,mirjalili2013kinetics} we define a simulation cell that resembles the interfacial structure of Ti--Al systems, as shown in \cref{fig:heat_treatment}.
% The simulation cell is defined as shown in \cref{fig:heat_treatment}.
The system has three layers, from top to bottom, there is a layer of pure Ti, a layer of TiAl\textsubscript{3} grains, and a layer of pure Al. The top \SI{1}{\nano\meter} layer of Ti and the bottom \SI{1}{\nano\meter} layer of Al are fixed in the xy-plane and cannot move freely. The z-axis movement is not restricted. A \SI{1}{\nano\meter}-wide margin of the TiAl\textsubscript{3} region is fixed in all three directions, preventing the TiAl\textsubscript{3} layer from floating freely in space.

\begin{figure}[htbp]
    \centering
    \includegraphics[width=0.6\textwidth]{figs/md_simulation_setup/heat_treatment_sys}
    \caption{The simulation cell for the heat treatment process. Besides the dimensions shown in the figure, the thicknesses of the Al, TiAl\textsubscript{3} and Ti layers are \SI{10}{\nano\meter}, \SI{4}{\nano\meter} and \SI{4}{\nano\meter} respectively. The TiAl\textsubscript{3} layer consists of $4$ randomly oriented grains, forming grain boundaries. The lattice orientations of the Al and Ti layers are also random. 
    % (The supporting information includes all lattice oritentations.)
    }
    \label{fig:heat_treatment}
\end{figure}

The interatomic interactions in the Ti--Al system are modeled using the embedded atom model (EAM) potential developed in \cite{zope2003interatomic} by fitting to both experimental and \textit{ab initio} data. This potential model has proven to be effective in various investigations and have been successfully tested for many mechanical and structural-energy properties of Ti-Al systems.\cite{poletaev2021self,wu2016molecular,tang2014molecular,gao2023design} A recent comprehensive study has demonstrated that after over two decades since its creation, this potential model remains a highly effective representation of the physical properties of the Ti--Al systems, even compared to newer models.\cite{pei2021systematic,kim2016atomistic}

The system is first equilibrated at \SI{300}{\kelvin} for \SI{1}{\nano\meter}. Then, it is heated up to \SI{895.15}{\kelvin} in \SI{300}{\pico\second}, corresponding to a heating rate of approximately \SI{2}{\kelvin/\pico\second}. Next, the system is kept at \SI{895.15}{\kelvin} for \SI{30}{\nano\second}. \SI{895.15}{\kelvin} is a typical temperature used in \cite{mirjalili2013kinetics}. Besides, multiple experimental studies show a temperature range of \SI{825}{\kelvin} to \SI{1220}{\kelvin} for Ti/Al system heat treatment.\cite{bizot2020reactivity,mirjalili2013kinetics,xu2006growth,sun2011multilayered,lazurenko2016explosively}

\subsection{Markov State Modeling methods}
\subsubsection{Modeling dynamics with graph representation of atoms}
To describe the state transition of a dynamical system, there exists a state space $\Omega$ that contains all dynamical variables needed, and $\bmx_t \in \Omega$ is a dynamical process. To describe the dynamics of an atom in a MD simulation, we can think of $\Omega$ as the space of local configurations of the atom. Such local configurations can often be encoded by the species of the target atom, the species of the surrounding atoms, distances (or bond length) between target atom and its neighbors, and sometimes, even the angular components describing the angles formed by atoms with target atom being the vertex. For example, \cref{asec:atomic_environment} describes how a local configuration of an atom can be a represented by a weighted colored graph.

% To build a Markov model that is robust in terms of the boundedness of the descretization error, Prinz \etal argued that $x(t)$ has to have three properties: 1. $x(t)$ is a Markov process 2. $x(t)$ is ergodic 3. $x(t)$ is reversible.

Then, the dynamics of the target atom can be described as a Markov process:\cite{xie2019graph}
\begin{equation}
    \bmx_{t+\tau} = \mathbf{F}(\bmx_t)
\end{equation}
where $\mathbf{F}$ is the dynamical system, which is non-linear in $\bmx$. Many recent studies have applied Koopman theory to build data-driven dynamical models. In short, Koopman theory states that there is a linearizing coordinate $\bmchi$ such that 
\begin{equation}
    \bmchi(\bmx_{t+\tau}) = \mathbf{K}(\tau)^\top\bmchi(\bmx_t)
\end{equation}
for some lag time $\tau$. It is shown that $\bfK$ has eigenvalues $1 \equiv \lambda_0 > \lambda_1 > \lambda_2 > \cdots > -1$.\cite{bowman2013introduction} The associated eigenvectors $\mathbf{v}_i$ can be viewed as principal modes of the dynamics, and the principal modes with the largest eigenvalues are the primary modes of probability flow between the system's substates \cite{bowman2013introduction}. Notably, $\mathbf{v}_0$ corresponds to the stationary distribution of the states, because raising $\lambda_0$ to any power does not change the weight of the process.\cite{prinz2011markov} For $i \neq 0$, the relaxation timescale of the $i$th slowest process is computed as
\begin{equation}
    t_i(\tau) = - \frac{\tau}{\ln \lvert \lambda_i(\tau) \rvert}
\end{equation}

% where the Koopman operator $\bfK$ is defined on the function space of $\bmchi$. 

% The VAMPnets framework learns $\bmchi$ from MD trajectories. The approximation of the $i$th eigenfunctions is 
% \begin{equation}
%     \hat{\psi}_i (\bmx) = \sum_j r_{ij} \bmchi_j (\bmx)
% \end{equation}
% {
%     \color{red}
    % what is $r_{ij}$?

    % eigen vectors?

    % decaying eigen values?

    % timescale calculation?
% }

In this work we use GDyNets developed in \cite{xie2019graph} to learn the Markov state model of the dynamics of the Ti-Al binary system. GDyNets is an implementation of the VAMP algorithm in TensorFlow dedicated to study atomistic level dynamics. The structure of the network and hyperparameter settings are described in \cref{asec:neural_network_architecture}.


\subsubsection{Sampling strategy}
For the training, testing and validation purpose of the MSM, we sample the interfacial dynamics of the Ti--Al binary system using a sandwich structure as illustrated in Figure \cref{fig:config}. The sandwich structure is composed of a substrate of either Ti or Al, with the top and bottom four layers of atoms being made of Ti in Figure \cref{fig:Ti-based-config} or Al in Figure \cref{fig:Al-based-config}. The middle of the sandwich structure is comprised of a mixture of Ti and Al atoms, with a total of $7500$ atoms.

\begin{figure}[htbp]
    \sidesubfloat[]{
        \includegraphics[height=8\baselineskip]{figs/sandwich_structure/sandwich_struct}\label{fig:sandwich_struct}
        }
    \sidesubfloat[]{
        \includegraphics[height=8\baselineskip]{figs/sandwich_structure/Ti-based}\label{fig:Ti-based-config}
        }
    \sidesubfloat[]{
        \includegraphics[height=8\baselineskip]{figs/sandwich_structure/Al-based}\label{fig:Al-based-config}
        }
    \caption{Schematics and examples of the sandwich configuration of the Ti-Al binary system. (a) Schematic of the sandwich structure; (b) The Ti-substrated configuration; (c) The Al-substrated configuration. The blue atoms represent the Ti atoms and the red atoms represent the Al atoms. A mixture of both species can be found in the middle of the sandwiches. In both configurations, the Ti-Al mixture consists of $7500$ atoms. The compositions of the mixtures are both $\mathrm{Ti}10\%\mathrm{Al}90\%$.}
    \label{fig:config}
\end{figure}

In our model, the substrates are treated as rigid bodies, while the $7500$-atom mixture is only subjected to thermostatting. Although in real experiments, no part of the system is truly rigid, this study places emphasis on the timescales of the transitions between dynamic states. The assumption of rigidity for the substrates is reasonable as the time required for the atoms in the solid substrate to transition into the liquid mixture is likely much longer than the other transitions occurring within the liquid mixture. Furthermore, modeling the substrates as rigid simplifies the system and ensures that it converges to a reversible dynamic equilibrium at a given temperature.

We focus on three key parameters of a sandwich structure: 1. the atomic species of the substrate $S = \{\mathrm{Ti}, \mathrm{Al}\}$, 2. the composition of the mixture $C = \{\mathrm{Ti}10\%\mathrm{Al}90\%, \mathrm{Ti}5\%\mathrm{Al}95\%, \mathrm{Ti}0\%\mathrm{Al}100\%\}$, and 3. the temperature $T = \{1033.5\mathrm{K}, 983.5\mathrm{K}\}$. Same as the temperature we set for the previous MD simulations, these two temperatures both fall in the range of typical Ti/Al system heat treatment. For an instance of the dynamical systems $s$, there is $s \in S \times C \times T$.

In the simulations, we use an overdamped Langevin thermostat, as suggested by \cite{prinz2011markov}. This ensures the dynamic process is Markovian, ergodic, and reversible, and thus satisfy the assumption of the theory of MSM developed in \cite{prinz2011markov}. For each instance of the dynamic system, we sample 5000 frames of the dynamics, with the time interval between every two frames being \SI{0.25}{\pico\second}. 2500 frames of the data are used for training, and 2000 for testing and 500 for validation.

% {
%     \color{red}
%     talk about number of frames in the dataset.

%     talk about the thermostat (overdamped langevin.) and so it satisfies detailed balance.
% }